(1) 均匀各向同性薄弹性板厚 \(h\)，面质量密度 \(\rho h\)，杨氏模量 \(E>0\)，Poisson比 \(-1<\nu<1/2\)。四边分别 \(x=0,a\) 与 \(y=0,b\) 均简支。采用Kirchhoff–Love小挠度理论，给定弯曲势能
\[
V=\frac D2\int_\Omega[w_{xx}^2+w_{yy}^2+2\nu w_{xx}w_{yy}+2(1-\nu)w_{xy}^2]\,dA,\quad D=\frac{Eh^3}{12(1-\nu^2)}.
\]
由变分得到内域场方程及简支边界条件。
(2) 求正交正弦本征模态与完整频率 \(\omega_{mn}\)，说明无转动和平移零频的理由。对方形板 \(a=b\) 讨论 \((m,n)\leftrightarrow(n,m)\) 的简并，并解释连续旋转对称为何并不存在。
(3) 在无限大同类薄板上求平面弯曲波色散 \(\omega(k)\)、相速度与群速度；核查长波群速，以及短波下忽略横向剪切和旋转惯量的适用界限。
(4) 由内域场方程与弯曲势能独立证明闭合简支板的总能量守恒；说明 \(-1<\nu<1/2\) 下势能正定的矩阵判据，并写出单一模态在任意时刻的动能与弯曲势能互换关系。
